\documentclass[12pt,a4paper]{article}

% =========================================
% PAGE SETUP
% =========================================
\usepackage[
    a4paper,
    margin=25.4mm
]{geometry}


% =========================================
% BASIC PACKAGES
% =========================================
\usepackage{graphicx}       % Images
\usepackage{amsmath}        % Mathematics
\usepackage{amssymb}        % Mathematical symbols
\usepackage{booktabs}       % Professional tables
\usepackage{float}          % [H] figure/table positioning
\usepackage{caption}        % Caption customization
\usepackage{hyperref}       % Hyperlinks and references
\usepackage{ragged2e}       % Text alignment
\usepackage{setspace}       % Line spacing
\usepackage{hyperref}       % Clickable links and references
% =========================================
% APA REFERENCING
% =========================================

\usepackage[
    style=apa,
    backend=biber
]{biblatex}

\addbibresource{references.bib}


% =========================================
% HYPERLINK SETTINGS
% =========================================
\hypersetup{
    colorlinks=true,
    linkcolor=black,
    urlcolor=black,
    citecolor=black,
    pdftitle={Critical Research Review Article}
}

% Table and figure caption formatting
\DeclareCaptionFont{captionnine}{\fontsize{9pt}{11pt}\selectfont}

\captionsetup[table]{
    font=captionnine,
    justification=raggedright,
    singlelinecheck=false,
    labelsep=colon
}

\captionsetup[figure]{
    font=captionnine,
    justification=centering,
    singlelinecheck=false,
    labelsep=colon
}


% =========================================
% DOCUMENT
% =========================================
\begin{document}

\begin{center}

{\fontsize{18}{22}\selectfont\textbf{CRITICAL RESEARCH REVIEW ARTICLE}}

\vspace{0.5cm}

{\fontsize{18}{22}\selectfont\textbf{A Comparative Review of Advanced Generative Adversarial Networks (GAN) Architectures}}

\vspace{1cm}

{\fontsize{12}{14}\selectfont\textbf{Group Members}}

\vspace{0.3cm}

{\fontsize{12}{14}\selectfont
1. Student Name -- Student ID\\
2. Student Name -- Student ID\\
3. Student Name -- Student ID\\
4. Student Name -- Student ID
}

\vspace{0.8cm}

{\fontsize{12}{14}\selectfont
\textbf{Programme:} [Programme Name]\\
\textbf{Course:} [Course Code and Course Name]\\
\textbf{Lecturer:} [Lecturer Name]\\
\textbf{Semester:} [Semester]\\
\textbf{Submission Date:} [Date]
}

\end{center}


% =========================================
% ABSTRACT
% =========================================

\vspace{0.5cm}

\noindent
\textbf{Abstract}

\vspace{0.2cm}

\begin{justify}

Write a concise summary of the entire review article. Write a concise summary of the entire review article. Write a concise summary of the entire review article. Write a concise summary of the entire review article.

Write a concise summary of the entire review article.Write a concise summary of the entire review article.Write a concise summary of the entire review article.






The abstract should include:

\begin{itemize}
    \item Background/context of the selected research topic
    \item Purpose of the review
    \item General approach used to select and analyse the research articles
    \item Major findings from the reviewed literature
    \item Key research gaps identified
    \item Overall conclusion
\end{itemize}

\textbf{Recommended length: 150--250 words.}

\vspace{0.3cm}

\textbf{Keywords:} Keyword 1; Keyword 2; Keyword 3; Keyword 4; Keyword 5

\end{justify}




% =========================================
% TABLE OF CONTENTS
% =========================================

\newpage

\tableofcontents

\newpage



% =========================================
% MAIN CONTENT
% =========================================

\begin{justify}

% =========================================
% 1. INTRODUCTION
% =========================================

\section{Introduction}

\subsection{Background}

Introduce the selected research topic and explain why it is important.

This is how you add inline references in the document given that all your references have been copied in the references.bib.
Recent studies have demonstrated the effectiveness of artificial
intelligence techniques \parencite{smith2024}.

Follow this example to include more than one references.Several studies have investigated these approaches
\parencite{smith2024,lee2023,ahmad2022}.

\vspace{0.5cm}
Discuss:

\begin{itemize}
    \item What is the research topic?
    \item What problem does it address?
    \item Why is the problem important?
    \item Where is the topic currently being applied?
    \item Why is there a need to study/review this topic?
\end{itemize}

Students should support important statements with appropriate research references.


\subsection{Problem Statement}

Clearly describe the problem or issue that motivates this literature review.

\vspace{0.5cm}
Explain:

\begin{itemize}
    \item What is currently known?
    \item What problems exist?
    \item What limitations exist in current approaches?
    \item Why is a review of existing research necessary?
\end{itemize}


\subsection{Objectives of the Review}

State the objectives of the review article.

\vspace{0.5cm}
For example:

\begin{enumerate}
    \item To review existing research related to [topic].
    \item To compare the approaches used by different researchers.
    \item To analyse the findings and limitations of existing studies.
    \item To identify research gaps and potential future research directions.
\end{enumerate}


\subsection{Research Questions}

Develop 2--4 research questions that guide the review.

\vspace{0.2cm}

For Example:


\textbf{RQ1:} What approaches have been used to address [research problem]?

\vspace{0.2cm}

\textbf{RQ2:} What are the major findings reported in existing studies?

\vspace{0.2cm}

\textbf{RQ3:} What are the limitations of the existing approaches?

\vspace{0.2cm}

\textbf{RQ4:} What research gaps and future opportunities can be identified?


% =========================================
% 2. REVIEW METHODOLOGY
% =========================================

\section{Review Methodology}

Explain how the group conducted the literature review.


\subsection{Literature Search Strategy}

Describe:

\begin{itemize}
    \item Databases/search engines used
    \item Search keywords
    \item Search strings
    \item Publication period
    \item Type of publications considered\newline
\end{itemize}


Example databases:

\begin{itemize}
    \item Scopus
    \item Web of Science
    \item IEEE Xplore
    \item ScienceDirect
    \item SpringerLink
    \item ACM Digital Library
    \item Google Scholar
\end{itemize}


\subsection{Inclusion and Exclusion Criteria}

Describe the criteria used to determine which research articles were included in or excluded from the review.

\textbf{Inclusion criteria may include:}

\begin{itemize}
    \item Articles directly related to the selected research topic
    \item Peer-reviewed research articles
    \item Articles published within the selected publication period
    \item Articles indexed in the selected academic databases
    \item Articles available in full text
\end{itemize}

\textbf{Exclusion criteria may include:}

\begin{itemize}
    \item Articles outside the scope of the research topic
    \item Duplicate publications
    \item Non-peer-reviewed publications
    \item Articles outside the selected publication period
    \item Articles without sufficient methodological or research information
\end{itemize}

\end{justify}

% =========================================
% SAMPLE TABLE
% =========================================

\begin{table}[H]
\caption{Inclusion and exclusion criteria used for the literature review}
\label{tab:inclusion-exclusion}

\fontsize{10pt}{12pt}\selectfont

\begin{tabular}{|p{4cm}|p{5cm}|p{5cm}|}
\hline
\centering\textbf{Criteria} &
\centering\textbf{Inclusion} &
\centering\textbf{Exclusion}
\tabularnewline
\hline

\centering Publication type &
\centering Peer-reviewed research articles &
\centering Non-peer-reviewed articles
\tabularnewline
\hline

\centering Publication period &
\centering [e.g., 2021--2026] &
\centering Articles outside selected period
\tabularnewline
\hline

\centering Language &
\centering English &
\centering Other languages
\tabularnewline
\hline

\centering Topic &
\centering Directly related to selected topic &
\centering Not directly related
\tabularnewline
\hline

\centering Full text &
\centering Full-text available &
\centering Full text unavailable
\tabularnewline
\hline

\end{tabular}

\end{table}

This is a sample table. Tables should be placed in the main text near to the first time they are cited, with the caption on top and left aligned. The font of the table caption is size 9 and the content in the table is size 10.



% =========================================
% SAMPLE FIGURE
% =========================================

\begin{figure}[H]
    \centering

    \fbox{%
        \includegraphics[width=0.80\textwidth]{figures/Prisma.jpg}%
    }

    \caption{PRISMA flow diagram of the study selection process.}
    \label{fig:prisma}
\end{figure}


This is a figure. All figures should be centred with figure caption following the sequence Figure 1, Figure 2, and so on. The caption should also be centred. The caption for figures should be size 9.



\subsection{Article Selection}

Each student must review \textbf{at least 3 research articles}.

Therefore:

\begin{itemize}
    \item \textbf{Minimum articles per student = 3}
    \item \textbf{Minimum articles per group = 12}
\end{itemize}

The group may use more than 12 articles if necessary.

Explain how the final articles were selected for the review.




% =========================================
% 3. LITERATURE REVIEW
% =========================================

\section{Literature Review}

This is the \textbf{main section of the review article}.

Do \textbf{NOT} simply write:

\begin{itemize}
    \item Article 1 says \ldots
    \item Article 2 says \ldots
    \item Article 3 says \ldots
\end{itemize}

Instead, organise the literature according to \textbf{themes, methods, technologies, problems, findings, or research approaches}.


\subsection{Theme 1: High Quality Synthetic Data Generation}

% Discuss research articles related to the first major theme.

% Compare the studies in terms of:

% \begin{itemize}
%     \item Research objectives
%     \item Methodology
%     \item Dataset/sample
%     \item Tools/technologies
%     \item Experimental setup
%     \item Results
%     \item Advantages
%     \item Limitations
% \end{itemize}

% The discussion should compare multiple studies rather than describe each article independently.

Below are the proposed advanced GAN architectures that are aimed to produce better synthetic data:
\begin{itemize}
    \item Dual Correlation Generative Adversarial Nets Based Virtual Sample Generation Method (DCRGAN-VSG)
    \item Parallelly Tempered Generative Adversarial Networks (PTGAN)
\end{itemize}

\subsubsection{Research Objectives}
\textbf{DCRGAN:} To address the small-sample problem in industrial soft sensing by generating additional virtual samples. The authors aim to improve conventional GAN-based sample generation by preserving the relationship between process variables and quality variables, while also selecting generated samples that fill sparse regions of the original dataset. \textbf{The overall goal is to create a more representative and balanced training dataset and thereby improve soft-sensor prediction accuracy.} \\
\\ \textbf{PTGAN:} To address the training instability of GANs, particularly when learning highly multimodal data distributions where problems such as mode collapse can occur. The authors propose a parallel tempering framework that uses multiple tempered distributions to reduce gradient variance and stabilize GAN training. \textbf{The overall goal is to enable GANs to learn complex target distributions more reliably and generate higher-quality synthetic data.}

\subsubsection{Methodology}
\textbf{DCRGAN}
\begin{enumerate}
    \item \textbf{Data Preparation:} The dataset is divided into training and testing sets.

    \item \textbf{DCRGAN Training:} Two generators and two discriminators are trained using a shared latent space. WGAN-GP generates process variables, while an improved Conditional GAN generates the corresponding quality variables. MSE loss is added to preserve the input-output relationship.

    \item \textbf{Virtual Sample Generation and Selection:} The trained DCRGAN generates paired virtual samples. The Pielou index evaluates their distribution, and only samples that improve dataset diversity and fill sparse regions are retained. This process continues until the required number of virtual samples is obtained.
\end{enumerate}
\textbf{PTGAN}
\begin{enumerate}
    \item \textbf{Constructing Intermediate Distributions:} PTGAN creates smoother intermediate distributions by interpolating between real samples using a temperature parameter.
    
    \item \textbf{Modifying the Generator:} The generator receives both the random noise and the temperature, allowing it to learn multiple distributions with a single model.
    
    \item \textbf{Parallel Adversarial Training:} The generator and critic are trained simultaneously across different temperature levels. The smoother distributions provide more stable gradients and reduce training instability.
    
    \item \textbf{Coherency Penalty:} A coherency penalty encourages the critic to learn consistently across temperature levels, allowing knowledge to be shared between them.
    
    \item \textbf{Final Data Generation:} After training, the generator is used to produce synthetic data that approximates the original data distribution while reducing mode collapse.
\end{enumerate}

\subsubsection{Dataset/Sample used}
\textbf{DCRGAN:} The DCRGAN-VSG method was evaluated using two types of datasets. 
\begin{itemize}
    \item First, a three-dimensional benchmark function was used as a simulation dataset to test whether the method could generate virtual samples with an appropriate distribution.
    \item Second, the method was tested on an industrial purified terephthalic acid (PTA) production dataset, where process variables and quality variables were used for soft-sensor modeling.
\end{itemize}
\textbf{PTGAN:} PTGAN was evaluated on both image and tabular datasets.
\begin{itemize}
    \item For image generation, it used CIFAR-10, BloodMNIST, and CelebA-HQ, which provide complex and multimodal data distributions.
    \item For tabular-data generation, it used the Adult, Credit Card Default, and Law School Admission datasets for supervised-learning tasks.
\end{itemize}

\subsubsection{Tools/Technologies used}
\textbf{DCRGAN:} DCRGAN combines several GAN-based techniques and mathematical methods to generate and select useful virtual samples:
\begin{itemize}
    \item WGAN-GP is used to generate process-variable samples.
    \item Conditional GAN (RC-CGAN) generates the corresponding quality variables.
    \item MSE loss is incorporated to preserve the relationship between process and quality variables.
    \item Pielou index is used to assess sample distribution and select samples that fill sparse regions.
\end{itemize}
\textbf{PTGAN:} PTGAN combines GAN-based generation, parallel tempering, and additional training techniques to improve gradient stability and synthetic data quality:
\begin{itemize}
    \item Parallel Tempering – creates smoother intermediate distributions to improve training stability.
    \item Temperature Parameter – allows the generator to learn across different distribution levels.
    \item Coherency Penalty – encourages consistent learning across different temperature levels.
\end{itemize}

\subsubsection{Experimental Setup}
\textbf{DCRGAN:} The DCRGAN was tested using a three-dimensional benchmark function and a real purified terephthalic acid (PTA) industrial dataset. They compared DCRGAN-VSG with several existing generative methods, including WGAN-GP, VAE, VAE-GAN, DDPM, and TTD. The generated virtual samples were then used to train soft-sensor models, and the prediction performance was compared to evaluate whether the proposed method improved soft sensing. \\
\\ \textbf{PTGAN:} PTGAN was evaluated on image and tabular datasets. For image generation, the experiments used CIFAR-10, BloodMNIST, and CelebA-HQ, with Inception Score (IS) and Fréchet Inception Distance (FID) used to measure generation quality. For tabular data, the researchers used the Adult, Credit Card Default, and Law School Admission datasets and evaluated the generated data using downstream random forest, SVM, and logistic regression models. The results were averaged over 10 independent runs to compare PTGAN with existing GAN training approaches.

\subsubsection{Results}
\textbf{DCRGAN:}
\begin{itemize}
    \item 70.43\% reduction in sensor estimation error on the 3D benchmark dataset.
    \item 39.33\% reduction in estimation error on the PTA industrial dataset.
    \item Achieved R² = 0.883 on the PTA dataset.
    \item Outperformed VAE, WGAN-GP, and DDPM in soft-sensing prediction performance.
\end{itemize}
\textbf{PTGAN:}
\begin{itemize}
    \item Achieved an FID of 41.99 on the BloodMNIST dataset, indicating improved image-generation quality.
    \item Reduced gradient variance, resulting in more stable GAN training through parallel tempering.
    \item Reduced demographic bias on the Law School dataset, achieving a statistical parity (SP) of 0.054.
    \item Demonstrated improved sample diversity and training stability compared with conventional GAN approaches.
\end{itemize}
\subsubsection{Advantages}
\textbf{DCRGAN:} Accurately learns the connection between measured inputs (process variables) and difficult-to-measure outputs (product quality) even when training on very small datasets. Its filtering step (Pielou index) ensures newly generated data actively fills empty gaps in the dataset instead of repeating what is already there. \\
\\ \textbf{PTGAN:} Solves unstable training and wild updates by training across blended, smoothed versions of the dataset simultaneously, which also allows it to generate fair and unbiased data across demographic groups.
\subsubsection{Limitations}
\textbf{DCRGAN:} Creating too many fake samples can backfire; generating more than roughly 420 samples distorts the original dataset’s shape and degrades overall model accuracy. \\
\\ \textbf{PTGAN:} Blending samples introduces slight approximation errors, meaning the model can lose sharp details if the balance between smoothed data and real data is not tuned carefully


\subsection{Theme / Area 2: [Insert Theme]}

Discuss the second major theme.

Compare and synthesise findings from different researchers.

Identify:

\begin{itemize}
    \item Similarities
    \item Differences
    \item Strengths
    \item Weaknesses
    \item Contradictory findings
    \item Trends
\end{itemize}


\subsection{Theme / Area 3: [Insert Theme]}

Discuss the third major theme.

Use evidence from multiple research articles.


\subsection{Theme / Area 4: [Insert Theme, if applicable]}

Add additional themes where appropriate.

\textbf{Note:} The number of themes should depend on the selected research topic. Do not create unnecessary sections simply to increase the number of pages.


% =========================================
% 4. COMPARATIVE ANALYSIS
% =========================================

\section{Comparative Analysis of Existing Studies}

Provide a systematic comparison of the reviewed research articles.

A table such as the following is recommended.


% =========================================
% COMPARATIVE ANALYSIS TABLE
% =========================================

\begin{table}[H]

\caption{Comparative analysis of the reviewed research studies}
\label{tab:comparative-analysis}

\resizebox{\textwidth}{!}{%
\begin{tabular}{|c|c|c|c|c|c|c|}
\hline

\textbf{No.} &
\textbf{Author \& Year} &
\textbf{Research Objective} &
\textbf{Method/Technique} &
\textbf{Dataset/Sample} &
\textbf{Key Findings} &
\textbf{Limitations}
\\
\hline

1 &
[Author, Year] &
[Research objective] &
[Method/Technique] &
[Dataset/Sample] &
[Key findings] &
[Limitations]
\\
\hline

2 &
[Author, Year] &
[Research objective] &
[Method/Technique] &
[Dataset/Sample] &
[Key findings] &
[Limitations]
\\
\hline

3 &
[Author, Year] &
[Research objective] &
[Method/Technique] &
[Dataset/Sample] &
[Key findings] &
[Limitations]
\\
\hline

4 &
[Author, Year] &
[Research objective] &
[Method/Technique] &
[Dataset/Sample] &
[Key findings] &
[Limitations]
\\
\hline

\end{tabular}%
}

\end{table}

% =========================================
% 5. SYNTHESIS OF FINDINGS
% =========================================

\section{Synthesis of Findings}

This section synthesises the collective observed and methodological findings from the reviewed advanced Generative Adversarial Network (GAN) architectures, drawing direct cross-comparisons across diverse data modalities, downstream evaluation settings, and structural paradigms.

\subsection{What Do the Studies Agree On?}

Across the reviewed literature, several core observations consistently emerge:

\begin{itemize}
    \item \textbf{Vulnerability of Standard GAN Formulations:} All studies uniformly agree that vanilla, unconstrained GAN architectures suffer from severe training instability, gradient variance collapse, and mode collapse when applied to complex or skewed target distributions \parencite{sohn2026parallelly,li2024improving}. Specifically, standalone objective functions optimise toward predominant modes, failing to capture sparse regions, tail-end minority classes, or fine-grained structural outliers \parencite{li2023eid}.
    
    \item \textbf{Insufficiency of Distribution Matching for Downstream Utility:} Simply matching global feature distributions is insufficient for task-specific performance. Studies in industrial soft-sensing \parencite{zhu2025novel}, defect inspection \parencite{li2023eid}, and clinical decision support \parencite{bietsch2023electronic} agree that naive adversarial training overlooks critical localised input--output mappings and intra-instance dependencies unless guided by explicit supervisory constraints or specialised loss regularisers.
    
    \item \textbf{Generative Sample Limits and Risk of Over-Generation:} Synthetic data only helps a model up to a certain point. After that, more is actually worse. For instance, feeding a sensor model too many fake readings (over roughly 420 samples for DCRGAN) blurs out the real physical boundaries it needs to measure accurately \parencite{zhu2025novel}. Similarly, simply generating more medical records doesn't help find rare, high-risk patients; it just clutters the dataset with useless, common details \parencite{bietsch2023electronic}.
    
    \item \textbf{Failure of Conventional Sample-Level Interpolation (SMOTE):} Multiple studies benchmark against traditional interpolation methods such as SMOTE, consistently observing that linear sample blending produces blurred textures, duplicates boundary noise, and lacks the capacity to capture high-dimensional non-linear distributions \parencite{li2023eid,zhu2025novel}.
\end{itemize}

\subsection{What Do the Studies Disagree On?}

Despite consensus on foundational vulnerabilities, the studies demonstrate notable divergences in architectural philosophy, loss formulation, and validation standards:

\begin{itemize}
    \item \textbf{Multi-Network Complexity vs. Unified Parameterisation:} A prominent methodological divergence exists between multi-agent frameworks and unified parameterisation. Approaches such as SOMGAN \parencite{li2024improving} and DCRGAN \parencite{zhu2025novel} address mode coverage and dual correlations by deploying groups of discriminators or multi-generator pairs. Conversely, models like PTGAN \parencite{sohn2026parallelly} argue that training instability can be solved within a single generator-critic framework through parallel continuous distribution tempering, avoiding the parameter footprint and structural overhead of discrete discriminator banks.
    
    \item \textbf{Efficacy of Norm Penalties in Adversarial Objectives:} Studies diverge sharply on how regression or distance constraints should be formulated. In fine-scale defect synthesis, \textcite{li2023eid} demonstrate that global $L_1$ and $L_2$ norm penalties completely fail because normal image backgrounds dominate optimisation, necessitating localised bounding-box mask penalties. In contrast, for continuous physical variable generation, \textcite{zhu2025novel} successfully rely on global Mean Squared Error (MSE) constraints to regularise conditional outputs.
    
    \item \textbf{Applicability and Faithfulness of Standard Metrics:} Research focusing on natural image synthesis considers Fr\'echet Inception Distance (FID) and Inception Score (IS) the gold standards for distribution fidelity \parencite{li2024improving,sohn2026parallelly}. In sharp contrast, specialised visual inspection \parencite{li2023eid} and clinical healthcare studies \parencite{bietsch2023electronic} disagree with the utility of standard FID, noting that low FID scores frequently fail to reflect downstream detection efficacy or clinical validity in extremely imbalanced or tabular contexts.
\end{itemize}


\subsection{Which Approaches Appear to Be Most Effective?}

The comparative evidence indicates that the most effective architectures are those that augment adversarial objectives with structural regularization tailored to the data modality:

\begin{itemize}
    \item \textbf{Topologically and Spatially Constrained Adversarial Networks:} For high-dimensional mode coverage, SOMGAN's self-organizing map topology outperforms independent multi-discriminator ensembles (such as Dropout-GAN and D2GAN) by preserving neighbor affinities and partitioning the discriminative feature space without redundant parameter overhead \parencite{li2024improving}.
    
    \item \textbf{Continuous Distribution Tempering for Gradient Variance Reduction:} PTGAN handles complex data across images and tables by pairing a parallel tempering objective with a coherency penalty, significantly cutting down variance and stabilising training. Rather than getting stuck in flat regions where learning stalls, it uses convex interpolation to bridge the gaps between data modes. This prevents vanishing gradients and allows the model to achieve a near-optimal minimax convergence rate \parencite{sohn2026parallelly}.
    
    \item \textbf{Dual-Supervised and Targeted Objective Regularization:} For physical and industrial applications, DCRGAN-VSG \parencite{zhu2025novel} and EID-GAN \parencite{li2023eid} demonstrate that injecting explicit geometric or domain-specific constraints (e.g., Pielou diversity filtering, auxiliary MSE loss, and cropped outlier penalties) significantly outperforms unconstrained GANs, reducing estimation errors in industrial soft-sensing by up to 70\% and achieving robust outlier retention.
\end{itemize}

\subsection{What Trends Can Be Observed?}

Analysis of recent literature reveals the following emerging trajectories in advanced GAN research:

\begin{itemize}
    \item \textbf{Technologies:} Integration of hybrid modeling techniques, notably coupling generative adversarial mechanisms with Vision Transformers (e.g., BEiT for evaluation), diffusion distillation, and modular generative--discriminative pipelines \parencite{sohn2026parallelly,li2024improving}.
    
    \item \textbf{Methodologies:} A pronounced transition from unguided adversarial generation toward constrained generation incorporating domain physics, topological graph structures, spatial masking, and temperature-controlled intermediate state sampling \parencite{sohn2026parallelly,zuo2023sketched}.
    
    \item \textbf{Datasets:} Expansion from standard benchmark image sets (e.g., MNIST, CIFAR-10) toward highly specialized industrial sensory logs (PTA plants), clinical Electronic Health Records (EHR), multi-view weakly annotated databases, and micro-scale defect images exhibiting extreme class imbalance \parencite{zhu2025novel,bietsch2023electronic,zhang2023inaccurate}.
    
    \item \textbf{Algorithms:} Replacement of standard Wasserstein weight clipping with advanced gradient stabilization frameworks, including coherency penalties, spectral normalization, and multi-agent competitive-cooperative loss functions \parencite{sohn2026parallelly,li2024improving}.
    
    \item \textbf{Applications:} Deployment in high-stakes operational environments, including industrial process soft-sensing \parencite{zhu2025novel}, automated manufacturing defect detection \parencite{li2023eid}, cost-sensitive facial biometric security \parencite{chen2021generative}, and military situation symbol standardization \parencite{zuo2023sketched}.
    
    \item \textbf{Research Directions:} A clear evolution toward \textit{Trustworthy and Pragmatic AI}, prioritizing algorithmic fairness (e.g., demographic parity via latent space interpolation), differential privacy guarantees in tabular medical synthesis, and task-specific downstream utility rather than purely aesthetic sample generation \parencite{sohn2026parallelly,bietsch2023electronic}.
\end{itemize}

% =========================================
% 6. RESEARCH GAPS
% =========================================

\section{Research Gaps}

Identify gaps that remain in the existing literature.

Possible research gaps include:

\begin{itemize}
    \item Lack of suitable datasets
    \item Small sample sizes
    \item Limited geographical coverage
    \item Poor generalisation
    \item Limited experimental validation
    \item Lack of comparison between methods
    \item High computational requirements
    \item Lack of real-world testing
    \item Security/privacy concerns
    \item Lack of explainability
    \item Outdated technologies
    \item Limited integration with other technologies
\end{itemize}

For each identified gap, provide evidence from the reviewed literature.


% =========================================
% 7. FUTURE RESEARCH DIRECTIONS
% =========================================

\section{Future Research Directions}

Based on the identified research gaps, propose potential directions for future research.

For example:

\begin{enumerate}
    \item Development of improved algorithms
    \item Use of larger datasets
    \item Integration of emerging technologies
    \item Development of real-time systems
    \item Improved accuracy/performance
    \item Cross-domain applications
    \item Improved security and privacy
    \item Explainable AI approaches
    \item Real-world deployment and validation
\end{enumerate}

The proposed directions should logically follow from the research gaps identified in Section~\ref{sec:research-gaps}.


% =========================================
% 8. DISCUSSION
% =========================================

\section{Discussion}

Provide an overall critical discussion of the literature.

The discussion should answer:

\begin{itemize}
    \item What has the research community achieved so far?
    \item Which approaches are most promising?
    \item What are the major challenges?
    \item Are the current approaches sufficient?
    \item What limitations remain?
    \item How has the research area evolved?
    \item What important issues remain unresolved?
\end{itemize}

This section should demonstrate \textbf{critical thinking}, not merely summarisation.


% =========================================
% 9. CONCLUSION
% =========================================

\section{Conclusion}

Summarise the major findings of the review.

The conclusion should:

\begin{itemize}
    \item Restate the purpose of the review
    \item Summarise the major findings
    \item Highlight important research gaps
    \item State the overall conclusion
    \item Briefly indicate future research opportunities
\end{itemize}

Do not introduce completely new information or references in the conclusion.






% =========================================
% REFERENCES
% =========================================

\newpage

\phantomsection
\addcontentsline{toc}{section}{References}

\printbibliography[
    title={References}
]





% =========================================
% APPENDIX A - INDIVIDUAL ARTICLE REVIEW
% =========================================

\newpage

\phantomsection
\addcontentsline{toc}{section}{Appendix A -- Individual Article Review}

\begin{flushleft}
    {\fontsize{14pt}{17pt}\selectfont
    \textbf{APPENDIX A -- INDIVIDUAL ARTICLE REVIEW}}
\end{flushleft}

\vspace{0.5cm}

Each student should complete the following table for their
\textbf{minimum 3 research articles} before the group combines the findings.

\vspace{0.5cm}

\textbf{Student:} [Name / Student ID]


\begin{table}[H]

\caption{Individual article review}
\label{tab:individual-article-review}

\fontsize{10pt}{12pt}\selectfont

\begin{tabular}{|p{1.5cm}|p{2.1cm}|p{1.8cm}|p{2.2cm}|p{2.1cm}|p{2.2cm}|p{2.0cm}|p{2.0cm}|}
\hline

\centering\textbf{Article} &
\centering\textbf{Research Problem} &
\centering\textbf{Objective} &
\centering\textbf{Methodology} &
\centering\textbf{Dataset/Sample} &
\centering\textbf{Key Findings} &
\centering\textbf{Limitations} &
\centering\textbf{Research Gap}
\tabularnewline
\hline

&  &  &  &  &  &  & 
\tabularnewline
\hline

&  &  &  &  &  &  & 
\tabularnewline
\hline

&  &  &  &  &  &  & 
\tabularnewline
\hline

&  &  &  &  &  &  & 
\tabularnewline
\hline

\end{tabular}

\end{table}


This appendix provides evidence of each student's individual contribution to the group review.



% =========================================
% APPENDIX B - ARTICLE MATRIX
% =========================================

\newpage

\phantomsection
\addcontentsline{toc}{section}{Appendix B -- Article Matrix}

\begin{flushleft}
    {\fontsize{14pt}{17pt}\selectfont
    \textbf{APPENDIX B -- ARTICLE MATRIX}}
\end{flushleft}

\vspace{0.5cm}

The group should combine all individual article reviews into one matrix.


\vspace{0.5cm}








\begin{table}[H]

\caption{Group article matrix of reviewed research studies}
\label{tab:article-matrix}

\fontsize{10pt}{12pt}\selectfont

\resizebox{\textwidth}{!}{%
\begin{tabular}{|c|c|c|c|c|c|c|c|c|}
\hline

\textbf{Article} &
\textbf{Student} &
\textbf{Year} &
\textbf{Method} &
\textbf{Dataset} &
\textbf{Key Results} &
\textbf{Strength} &
\textbf{Limitation} &
\textbf{Research Gap}
\\
\hline

1 & Student 1 & & & & & & & 
\\
\hline

2 & Student 1 & & & & & & & 
\\
\hline

3 & Student 1 & & & & & & & 
\\
\hline

4 & Student 2 & & & & & & & 
\\
\hline

5 & Student 2 & & & & & & & 
\\
\hline

6 & Student 2 & & & & & & & 
\\
\hline

... & ... & & & & & & & 
\\
\hline

\end{tabular}%
}

\end{table}

\end{document}
